%%%PAGE 263 rot=0 legibility=good summary=圆周运动：等效重力场临界速度、水平面转弯（汽车/倾斜路面）、杆模型、绳不松弛临界与脱离角
%%%BLOCK id=263-01 ch=04 sec=竖直面圆周运动·等效重力场 kind=method alt=09 cont=prev y=0.00-0.31
%FIG p263_a 0.06 0.005 0.18 0.115
\notefig[0.25]{figures/p263_a.png}
$\downarrow$ $E=\dfrac{2mg}{q}$
\[
\left\{\begin{aligned}
&mg+qE+F=\dfrac{mv^2}{L}\\
&T\ge 0\qquad V_{\text{上}}\ge\sqrt{3gL}
\end{aligned}\right.
\]
% CHECK: 第一式写 F、条件写 T≥0，应为同一个量（绳的拉力）
\[ G'=3mg,\qquad g'=3g,\qquad V_{\text{下}}\ge\sqrt{15gL} \]
% CHECK: g'=3g 中“3”前有涂改痕迹

（图注：等效最高点）
%FIG p263_b 0.03 0.108 0.205 0.215
\notefig[0.3]{figures/p263_b.png}
\[ \sqrt2\,mg+F=\dfrac{mv^2}{L}\qquad F\ge 0 \]
\[ V\ge\sqrt{\sqrt2\,gL} \]
%FIG p263_c 0.09 0.22 0.22 0.315
\notefig[0.25]{figures/p263_c.png}
\[ mg\left(L+\dfrac{\sqrt2}{2}L\right)+qE\,\dfrac{\sqrt2}{2}L=\dfrac12m\left(V_0^2-V^2\right) \]
%%%END
%%%BLOCK id=263-02 ch=04 sec=水平面圆周运动·转弯 kind=method alt=02 cont=none y=0.32-0.59
\textbf{水平面圆周运动}\quad 必有侧\textsuperscript{向外}向滑动趋势
% CHECK: “向外”为行上方的插入小字，插入位置约在“侧”“向”之间，语义待核（可能想表达“必有向外侧滑的趋势”）

\textbf{汽车转弯}\quad 水平路面上：
%FIG p263_d 0.17 0.388 0.285 0.466
\notefig[0.25]{figures/p263_d.png}
\[ f=\dfrac{mv^2}{r}\le\mu mg\qquad v\le\sqrt{\mu gR} \]

\textbf{倾斜路面上}\quad 无侧滑趋势\quad $V_0$.
%FIG p263_e 0.335 0.485 0.495 0.59
\notefig[0.3]{figures/p263_e.png}
\begin{align*}
N\sin\alpha&=\dfrac{mv_0^2}{R}\\
N\cos\alpha&=mg
\end{align*}
\[ V_0=\sqrt{gR\tan\alpha}\qquad\text{与 }\mu\text{ 无关} \]
%%%END
%%%BLOCK id=263-03 ch=04 sec=竖直面圆周运动·杆模型 kind=method alt=none cont=none y=0.60-0.78
\textbf{杆模型}
%FIG p263_f 0.06 0.635 0.285 0.72
\notefig[0.3]{figures/p263_f.png}
以向下为正
\[ mg+F=\dfrac{mv^2}{R}\qquad F\text{向下} \]
\[ mg-(-F)=\dfrac{mv^2}{R}\qquad F\text{向上} \]
\[
\left\{\begin{aligned}
&0\le v<\sqrt{gR}\quad\text{推力}\quad v\uparrow\ \text{推力}\downarrow\\
&v>\sqrt{gR}\quad\ \ v\uparrow\ F_{\text{拉}}\downarrow
\end{aligned}\right.
\]
% CHECK: v>√(gR) 时 F拉 随 v 增大应增大，原稿箭头画的是 ↓
%FIG p263_g 0.70 0.615 0.885 0.72
\notefig[0.3]{figures/p263_g.png}
\[ F=\dfrac{mv^2}{R}-mg \]
%%%END
%%%BLOCK id=263-04 ch=04 sec=竖直面圆周运动·绳模型临界 kind=method alt=none cont=none y=0.80-0.92
\textbf{绳不松弛}\quad $N=0$ 临界.
%FIG p263_h 0.07 0.86 0.155 0.92
\notefig[0.25]{figures/p263_h.png}
\[ V_0\le\sqrt{2gL} \]
\[ V_0\ge\sqrt{5gL} \]
%%%END
%%%BLOCK id=263-05 ch=04 sec=竖直面圆周运动·脱离临界 kind=method alt=06 cont=none y=0.92-1.00
%FIG p263_i 0.09 0.93 0.285 1.0
\notefig[0.3]{figures/p263_i.png}
\[ V_0=\sqrt{\dfrac{7}{2}gR} \]
\begin{align*}
mg\sin\alpha&=\dfrac{mv^2}{R}\\
-mgR(1+\sin\alpha)&=\dfrac12m\left(gR\sin\alpha-\dfrac{7}{2}gR\right)
\end{align*}
% CHECK: 第二式 mg 后的 R 与左括号重叠，按能量关系取 R
\[ \Rightarrow\ \alpha=30^\circ \]
%%%END
%%%PAGEEND 263
